\chapter{Conclusiones}\label{cap:conclusiones}

\section{Respuesta a la pregunta}

En los condados de Estados Unidos, la mortalidad por cáncer se asocia, a igualdad del resto de
características, con:
\begin{itemize}
  \item la \textbf{incidencia}, que es el factor dominante: un condado en el tercer cuartil de
  incidencia tiene \res{efiqr-incidencia} muertes por \num{100000} habitantes más que uno en el
  primero;
  \item la \textbf{cobertura sanitaria sólo pública}, indicador de pobreza y peor acceso a la
  atención especializada (\res{efiqr-solopublico} muertes más entre el primer y el tercer cuartil);
  \item la \textbf{educación}: más adultos con título universitario, menos mortalidad
  (\res{efiqr-grado25-sur} muertes entre cuartiles en el Sur), con un gradiente mucho más débil en
  el Noreste; y más jóvenes cuyo nivel máximo es el bachillerato, más mortalidad;
  \item el \textbf{desempleo}, pero sólo en el Medio Oeste y el Noreste;
  \item la \textbf{región}: a igualdad de composición, el Sur conserva una mortalidad mayor que el
  Noreste y el Oeste, aunque el modelo explica el \res{pct-brecha-explicada}\,\% de la brecha bruta
  entre el Sur y el Oeste.
\end{itemize}

El modelo final, estimado por MCO con \res{n-final} condados, explica el \res{pct-r2-final}\,\% de
la variabilidad ($\bar R^2=\res{r2aj-final}$, $\mathrm{RMSE}=\res{rmse-final}$). Los errores son
heterocedásticos y dependen dentro de cada estado; la inferencia lo tiene en cuenta mediante
errores típicos robustos por conglomerados, y las conclusiones se mantienen en las
\res{n-especificaciones} especificaciones del análisis de sensibilidad y en el bootstrap por
conglomerados.

\section{Limitaciones}

\begin{itemize}
  \item \textbf{Diseño ecológico y observacional.} Las asociaciones describen condados, no
  personas, y no son efectos causales: un condado con más titulados universitarios difiere en
  muchas otras cosas no medidas (tabaquismo, obesidad, exposición laboral, acceso a cribados).
  \item \textbf{Variables omitidas.} No se dispone de prevalencia de tabaquismo, el principal factor
  de riesgo evitable de muerte por cáncer, ni de la distribución por tipos de cáncer o por edad.
  \item \textbf{Ausencia estructural de la incidencia.} Los condados de los estados cuyos registros
  no publican la incidencia quedan fuera del modelo principal; la imputación múltiple muestra que
  las conclusiones no cambian, pero bajo el supuesto MAR.
  \item \textbf{Selección de variables.} Los errores típicos se calculan condicionados al modelo
  seleccionado; la reestimación del modelo máximo y la concordancia entre estrategias de selección
  limitan, pero no eliminan, este problema.
\end{itemize}

\section{Trabajo futuro}

Los datos admiten una estimación más eficiente si se modela directamente el número de muertes con
un modelo de Poisson o binomial negativo con la población como exposición, lo que incorpora de forma
natural la heterocedasticidad debida al tamaño del condado; un modelo espacial permitiría además
modelar la dependencia entre condados vecinos de estados distintos.
